\section{Ablation Studies}
\label{sec:ablation}

We validate the key design choices of \method{} through two ablation experiments on Qwen2.5-7B. Table~\ref{tab:ablation} summarizes the results; training curves are provided in Appendix~\ref{sec:appendix-results}.

\begin{table}[t]
\centering
\small
\setlength{\tabcolsep}{4pt}
\begin{tabular}{lcccc}
\toprule
\textbf{Method} & \textbf{AIME} & \textbf{AIME} & \textbf{Olympiad} & \textbf{MATH} \\
& \textbf{2024} & \textbf{2025} & \textbf{Bench} & \textbf{-500} \\
\midrule
ORM & 10.8 & 10.8 & 46.3 & 80.2 \\
\midrule
Mult ($r^{\text{out}} \!\times\! r^{\text{proc}}$) & 12.1 & 10.2 & 46.7 & 80.2 \\
Fullnorm & 12.5 & 10.4 & 49.6 & 81.2 \\
\textbf{\method{}} & \textbf{15.8} & \textbf{13.1} & \textbf{51.3} & \textbf{82.3} \\
\bottomrule
\end{tabular}
\caption{Ablation results (\%, avg@4, Qwen2.5-7B). All methods report accuracy.}
\label{tab:ablation}
\end{table}

\paragraph{Correct-subset vs.\ full normalization.}
The \textbf{Fullnorm} variant normalizes $\aproc$ over all $G$ responses including incorrect ones. It performs slightly below \method{} across all benchmarks. The gap arises because including incorrect responses ($r^{\text{proc}} = 0$) in the normalization causes $\aproc$ to partially recapitulate the correct-vs-incorrect distinction already captured by $\aout$, diluting the fine-grained quality signal.

\paragraph{Multiplicative reward baseline.}
The \textbf{Mult} variant combines rewards via $r_i = r_i^{\text{out}} \times r_i^{\text{proc}}$ with standard GRPO normalization. Mult performs comparably to ORM and falls well below \method{}. In a single normalization pass, the outcome difference dominates in mixed-correctness groups, effectively suppressing the process quality information.

Both ablations confirm that \method{}'s gains require both \emph{decoupled normalization}, validated by Mult's inferior performance, and \emph{correct-subset normalization}, validated by Fullnorm's inferior performance, working in concert.
